\begin{figure}[htbp]
\centering
\begin{tikzpicture}[>=Stealth,every node/.style={font=\small,align=center},
block/.style={draw=ink,rounded corners=2pt,inner sep=6pt,text width=35mm},
wide/.style={draw=ink,inner sep=6pt,text width=47mm}]
\node[block] (d) {État arithmétique\\admissible\\\((x,y,c)\)};
\node[block,right=14mm of d] (e) {État raffiné\\\((10x-c,10y+c)\)};
\node[block,right=14mm of e] (j) {Bases d’étiquettes\\\(J\) unitaire};
\draw[<->] (d)--node[above,font=\scriptsize]{retenue} (e);
\draw[->] (e)--(j);
\node[wide,below=17mm of e] (w) {Transport bilatéral\\\(W^*W=I\)};
\draw[->] (j.south)|-(w.east);
\node[block,below left=17mm and 4mm of w] (a) {Lecture conjointe\\\(\frac{72G+R^*GR}{73}\)};
\node[block,below right=17mm and 4mm of w] (b) {Observable transporté\\\(WGW^*\)};
\draw[->] (w)--(a);\draw[->] (w)--(b);
\node[font=\small,below=8mm of a] {Invariant si \(GR=RG\)};
\node[font=\small,below=8mm of b] {Récupération au moyen de \(W^*\)};
\end{tikzpicture}
\caption{Composition d’opérations démontrées. Le raffinement conserve la somme normalisée et permet de récupérer l’étiquette. Son relèvement agit sur les amplitudes des états et conserve la norme. La lecture d’un observable et son transport complet conservent des fonctions différentes, spécifiées par les deux formules inférieures.}
\label{libro:fig:composicion}
\end{figure}
